Subtracting a reference orbit and dividing by its first variation removes
the affine part of a scalar transition. The remaining equation vanishes
to second order in the normalized phase variable. An integrable bound on that
equation gives uniform endpoint maps even when the unnormalized first
variation has no uniform bound.

\begin{lemma}[Normalized scalar transitions]\label{lem:scalar-normalization}
Let \(k\ge1\). Consider equations
\[
 \eta_r=N_\vartheta(r,\eta),\qquad L\le r\le R_\vartheta,
 \qquad N_\vartheta(r,0)=\partial_\eta N_\vartheta(r,0)=0,
\]
on \(|\eta|\le a\), where \(L\) and \(a>0\) are fixed.
Suppose that, uniformly in the parameters,
\[
 \|N_\vartheta(r,\cdot)\|_{C^{k+1}([-a,a])}\le g(r),
 \qquad g\in L^1([L,\infty)).
\]
Then there is a fixed neighborhood of zero on which all endpoint
maps have bounded \(C^{k+1}\) norms and positive first derivatives
bounded away from zero. Their inverses have the same bounds on a
fixed smaller neighborhood of zero.

Suppose also that \(R_\vartheta\to\infty\) and the kernels converge
in \(C^0([L,M];C^{k+1}_\eta)\) for every fixed \(M>L\).
Then the endpoint maps
and their inverses converge in \(C^{k+1}\) on smaller neighborhoods.
For a chosen parameter derivative \(\partial_\ell\), convergence
in \(C^k\) holds for the differentiated maps if the differentiated
kernels converge in \(C^0([L,M];C^k_\eta)\), have a common integrable
majorant, and
\[
 \|N_\vartheta(R_\vartheta,\cdot)
               \partial_\ell R_\vartheta\|_{C^k}\longrightarrow0.
\]
The convergence is joint in an additional compact parameter family
if these local convergences, \(R_\vartheta\to\infty\), and the
vanishing endpoint term are uniform in that family.
\end{lemma}
\begin{proof}
Taylor's formula gives \(|N_\vartheta(r,\eta)|\le g(r)|\eta|^2/2\).
Comparison with \(y'=g(r)y^2/2\) keeps all solutions from a common
enlarged initial interval inside \(|\eta|<a\). The first variation is
\[
 \exp\left(\int_L^r
          \partial_\eta N_\vartheta(t,\eta(t))\,dt\right)>0.
\]
It and its reciprocal are bounded. Induction in the triangular
variational equations and Gronwall's inequality bound all remaining
rows through order \(k+1\). The slope lower bound gives a common
image interval, and the inverse identity gives the inverse bounds.

Subtracting the integral equations and their variations gives convergence
on \([L,M]\); each remaining tail costs at most \(C\int_M^\infty g\).
Let the parameters tend to their limit and then \(M\to\infty\).
For the parameter derivative, the variation satisfies
\[
 Z_r=\partial_\eta N_\vartheta\,Z+\partial_\ell N_\vartheta,
 \qquad Z(L)=0.
\]
The endpoint derivative adds
\(N_\vartheta(R_\vartheta,\eta(R_\vartheta))
\partial_\ell R_\vartheta\). The same argument proves \(C^k\)
convergence under the endpoint assumption: differentiating a coefficient
\(N_{\eta^h}(r,\eta)\) adds
\(N_{\eta^{h+1}}(r,\eta)Z\), which requires the extra phase order.
Finally, \(C^{k+1}\) convergence of the endpoint maps \(A_\vartheta\) and
\[
 \partial_\ell A_\vartheta^{-1}
 =-\left(\frac{\partial_\ell A_\vartheta}{A_\vartheta'}\right)
     \circ A_\vartheta^{-1}
\]
give \(C^k\) convergence of the differentiated inverses. Each estimate
is uniform in the additional compact family under the stated assumptions.
\end{proof}

For a graph equation \(Y_r=F(r,Y)\), a reference solution \(Y_*\)
and its first variation \(j>0\), the substitution
\(\eta=(Y-Y_*)/j\) gives the exact kernel
\begin{equation}\label{tail:normalized-kernel}
 N(r,\eta)=j\eta^2\int_0^1(1-t)
                  F_{YY}(r,Y_*+tj\eta)\,dt.
\end{equation}
For \(h\ge2\), \(\partial_\eta^hN=j^{h-1}F_{Y^h}(r,Y_*+j\eta)\).
At fixed \(r,\eta\), a parameter derivative is therefore
\begin{equation}\label{deg:det-parameter-phase-derivative}
\begin{split}
 \partial_\ell\partial_\eta^hN
 =j^{h-1}\bigl\{&(h-1)(\partial_\ell\log j)F_{Y^h}
 +\partial_\ell^{\rm exp}F_{Y^h}\\
 &+F_{Y^{h+1}}(Y_{*,\ell}+j_\ell\eta)\bigr\}.
\end{split}
\end{equation}
Here \(\mathrm{exp}\) holds the graph coordinates fixed. The orders
zero and one follow by integrating in \(\eta\), since
\(N(r,0)=N_\eta(r,0)=0\). Thus \(C^3\)-valued parameter
derivatives require \(F_{Y^4}\) and \(N_{\eta^4}\).
Integrating the indicated parameter derivative gives the sequential
parameter secants used below. This requires neither a mixed parameter
derivative nor convergence of \(j\).

A quantitative form controls how close the normalized map is to the
identity when only an upper bound on the first variation is available.
For an actual graph family
write \(J_n=\partial_b^nY\), \(Q_n=J_n/J_1\), and
\(a_n=F_{Y^n}J_1^{n-1}\). Direct differentiation gives
\begin{equation}\label{tail:relative-jet-equations}
 Q_2'=a_2,\qquad Q_3'=3a_2Q_2+a_3,\qquad
 Q_4'=a_2(4Q_3+3Q_2^2)+6a_3Q_2+a_4.
\end{equation}
If, uniformly on an enlarged label interval,
\(Q_n(L)=O(\varepsilon)\) and
\(\int|a_n|\,dr=O(\varepsilon)\), \(2\le n\le4\),
these identities give \(Q_n=O(\varepsilon)\).
Since \(\partial_b\log J_1=Q_2\), the endpoint map in \(b-b_*\),
normalized by its exact value and slope at \(b_*\), and its inverse on a
smaller common interval, are \(\mathrm{id}+O_{C^3}(\varepsilon)\)
for sufficiently small \(\varepsilon\).
This conclusion is independent of the size of \(J_1(b_*)>0\).
